\section{Introduction}\label{sec:intro}

Data centres are among the fastest-growing electricity consumers. Efficiency gains kept their energy use nearly flat through the 2010s even as computing output grew several-fold~\cite{masanet2020recalibrating,shehabi2016united}. That buffer is now thinning, and information and communication technology accounts for an estimated 2.1--3.9\% of global greenhouse-gas emissions~\cite{freitag2021real}. Large machine-learning workloads have drawn particular attention~\cite{strubell2019energy,patterson2022carbon,dodge2022measuring}. The carbon intensity of grid electricity varies by a factor of ten or more across regions and hours, driven by the output of wind, solar and dispatchable plant~\cite{staffell2017measuring}. Operators can therefore cut emissions without new hardware by moving flexible work to cleaner places and times. Such \emph{carbon-aware computing}~\cite{radovanovic2023carbon,bashir2021enabling} is an operations problem: capacity is finite, deadlines bind, and both the grid and the job stream are uncertain.

Batch jobs such as model training, analytics, rendering and backups are a natural target, because they tolerate some delay. Previous work has shown the potential of temporal shifting (running later)~\cite{wiesner2021lets,goiri2012greenhadoop,hanafy2023carbonscaler} and of spatial shifting (running elsewhere)~\cite{liu2011greening,zheng2020mitigating}. It has also shown their limits. Sukprasert et al.~\cite{sukprasert2024limitations} found that practical reductions from spatio-temporal shifting fall well short of idealised bounds. Most methods are evaluated with a scheduler that reacts to the jobs it has already seen, using point forecasts of carbon intensity. Two questions have received less attention. First, how much is lost because an online scheduler does not know which jobs will arrive next? Second, does it pay to price the uncertainty of carbon-intensity forecasts into scheduling decisions?

This paper answers both questions with a new algorithm and a forecasting model, evaluated on three years of real regional grid data for Great Britain. It makes four contributions.
\begin{enumerate}
\item \textbf{Formulation.} We state multi-site, deadline-constrained, width-limited scheduling of deferrable bag-of-tasks jobs, including the emissions of data transfer, as a minimum-cost flow problem (Proposition~\ref{prop:flow}). Offline clairvoyant optima and every online planning step can therefore be computed exactly and quickly by network simplex. We also prove that myopic receding-horizon scheduling can be arbitrarily worse than the clairvoyant optimum even with perfect carbon forecasts (Proposition~\ref{prop:myopic}).
\item \textbf{Algorithm.} We propose CARMA (Carbon-aware Anticipatory Receding-horizon Multi-site Allocation). At every hour, CARMA adds \emph{ghost jobs} to its planning network. These represent expected arrivals, learned from the operator's history, so the planner reserves scarce low-carbon capacity for work that has not yet arrived. The planning problem stays a single min-cost flow and decisions take tens of milliseconds.
\item \textbf{Forecasting model.} We develop a probabilistic 1--24-hour carbon-intensity forecaster. It consists of direct gradient-boosted models with horizon-wise normalised conformal calibration and adaptive conformal inference, which holds coverage under the strong decarbonisation drift seen in 2024. Its calibrated quantiles let us test risk-adjusted carbon pricing inside the scheduler.
\item \textbf{Evidence.} Across 52 out-of-sample weeks, three load levels and two operating regimes, we compare CARMA with five implementable online baselines, clairvoyant and ablated variants, and the clairvoyant oracle. The results separate the value of anticipating demand from the value of better carbon forecasts, and show when each matters.
\end{enumerate}

The results carry an operational message. Once spatial shifting is allowed, the dominant loss comes from demand myopia, not from forecast error. Anticipating demand cuts the gap to the clairvoyant optimum by a factor of 2.2--3.0. By contrast, perfect carbon foresight or risk-averse carbon pricing adds little once demand is anticipated. Without migration, temporal shifting alone achieves much smaller reductions. The forecast-driven planners are then within 4--6\% of the optimum, and the remaining gap is due almost entirely to forecast error. The rest of the paper is organised as follows. Section~\ref{sec:related} reviews related work. Section~\ref{sec:methods} presents the model, the forecaster, CARMA and the experimental design. Section~\ref{sec:results} reports results, Section~\ref{sec:discussion} discusses implications and limitations, and Section~\ref{sec:conclusion} concludes.

\section{Related work}\label{sec:related}

\paragraph{Carbon-aware workload shifting} Early green data-centre schedulers matched deferrable jobs to on-site solar production, for example GreenSlot~\cite{goiri2011greenslot} and GreenHadoop~\cite{goiri2012greenhadoop}. Other work exploited electricity-price and carbon differences across sites through geographical load balancing~\cite{rao2010minimizing,liu2011greening} and dynamic right-sizing~\cite{lin2013dynamic}. With public carbon-intensity feeds, attention moved to grid emissions. Wiesner et al.~\cite{wiesner2021lets} quantified the reductions from temporal shifting in several regions, including Great Britain, and proposed forecast-based lowest-window scheduling. Google's carbon-intelligent compute system sets next-day virtual capacity curves from carbon and demand forecasts~\cite{radovanovic2023carbon}. CarbonScaler varies the allocation of elastic jobs with carbon intensity~\cite{hanafy2023carbonscaler}. Ecovisor exposes the energy system to applications~\cite{souza2023ecovisor}. Carbon Explorer weighs scheduling against investment in renewables and storage~\cite{acun2023carbon}. Load migration between data centres can absorb curtailed renewable generation~\cite{zheng2020mitigating}, and data-centre demand response has been surveyed in~\cite{wierman2014opportunities}. Sukprasert et al.~\cite{sukprasert2024limitations} bound the potential of spatio-temporal shifting across 123 regions and find that simple policies capture most of the achievable benefit. Our study complements theirs by identifying \emph{which} information limits online schedulers when capacity at clean sites is scarce.

\paragraph{Energy-aware scheduling in operations} Energy-efficient scheduling is a mature topic in manufacturing operations~\cite{gahm2016energy}, and green scheduling models have appeared in this journal, for example for parallel machines with job splitting~\cite{asadpour2022green} and for cloud resource allocation~\cite{sharma2024efficient}. Most such models are deterministic and solved offline, often with metaheuristics. We instead exploit exact network-flow structure~\cite{ahuja1993network,orlin1997polynomial} and embed it in an online lookahead policy. In the taxonomy of Powell~\cite{powell2019unified}, CARMA is a deterministic lookahead policy with a modified (anticipatory) model, in the spirit of model predictive control~\cite{rawlings2017model,mayne2000constrained}.

\paragraph{Carbon-intensity forecasting} Forecasts of grid carbon intensity use statistical decomposition~\cite{bokde2021short}, machine learning with many exogenous drivers~\cite{leerbeck2020short}, day-ahead regression for building demand response~\cite{lowry2018day}, and hierarchical models of the generation mix, as in CarbonCast~\cite{maji2022carboncast}. For Great Britain, NESO publishes regional forecasts through its Carbon Intensity API~\cite{nesoapi}. These studies focus on point accuracy. For decisions, calibrated uncertainty matters as well. Conformal prediction~\cite{vovk2005algorithmic,lei2018distribution,angelopoulos2023conformal} provides distribution-free intervals and has been extended to quantile regression~\cite{romano2019conformalized}, to multi-horizon~\cite{stankeviciute2021conformal} and dynamic time series~\cite{xu2023conformal}, and to arbitrary distribution shift through adaptive conformal inference~\cite{gibbs2021adaptive}. We use these tools to build a calibrated forecaster and to test, in a controlled way, whether risk-adjusted carbon costs improve scheduling. The question is motivated by the optimizer's curse~\cite{smith2006optimizer} and by the robust-optimisation view that point-estimate solutions can be fragile~\cite{bental2000robust,bertsimas2004price}.
